% !TEX root = ../main.tex
\begin{ouverture}[mole]{1}{Composition d'un système chimique}{Partie 1 · Constitution et transformations de la matière}
{Mole\enspace·\enspace Masse molaire\enspace·\enspace Volume molaire\enspace·\enspace Concentration\enspace·\enspace Dissolution et dilution}
\objectif{Déterminer la masse molaire d'une espèce à partir des masses molaires atomiques des éléments qui la composent.}
\objectif{Relier la masse d'un échantillon, sa quantité de matière et le nombre d'entités qu'il contient.}
\objectif{Utiliser le volume molaire d'un gaz pour déterminer une quantité de matière.}
\objectif{Déterminer la quantité de matière de chaque espèce d'un mélange à partir de sa composition.}
\objectif{Déterminer la quantité de matière d'un soluté à partir de sa concentration et du volume de solution.}
\prerequis{La mole et la constante d'Avogadro, la concentration en masse (classe de seconde).}
\prerequis{Préparation d'une solution par dissolution et par dilution (classe de seconde).}
\prerequis{Puissances de dix, conversions d'unités de volume et de masse.}
\end{ouverture}
\essentiel
\essSection{1}{La quantité de matière}
\begin{definition}[La mole]
Les entités chimiques (atomes, molécules, ions) sont trop petites et trop nombreuses pour être comptées une à une : on
les compte par paquets. Une \textbf{mole} est un paquet de $\num{6.02e23}$ entités ; la \textbf{constante d'Avogadro} vaut
$N_\mathrm{A}=\qty{6,02e23}{mol^{-1}}$. La \textbf{quantité de matière} $n$ d'un échantillon, en moles, est le nombre de paquets qu'il contient :
\[N=n\times N_\mathrm{A},\]
où $N$, sans unité, est le nombre d'entités de l'échantillon.
\end{definition}
\schema{01-relations}
\essSection{2}{Masse molaire et quantité de matière}
\begin{definition}[Masse molaire]
La \textbf{masse molaire atomique} d'un élément est la masse d'une mole de ses atomes ; elle est donnée par le tableau
périodique, en \unit{g.mol^{-1}}. La \textbf{masse molaire} $M$ d'une espèce chimique est la masse d'une mole de ses entités : c'est la
somme des masses molaires atomiques des atomes qui composent l'entité. La masse des électrons étant négligeable, un ion a
la même masse molaire que l'atome ou le groupe d'atomes dont il provient.
\end{definition}
\begin{propriete}[Relation entre masse et quantité de matière]
Pour un échantillon de corps pur de masse $m$ (en g) et de masse molaire $M$ (en \unit{g.mol^{-1}}) :
\[n=\frac{m}{M}.\]
Pour un liquide de volume $V$ et de masse volumique $\rho$, on calcule d'abord sa masse : $m=\rho\times V$.
\end{propriete}
\begin{savoirfaire}[Calculer une quantité de matière à partir d'une masse]
On cherche le nombre de molécules de saccharose, $\ce{C12H22O11}$, dans un morceau de sucre de \qty{10,0}{g}.
\begin{questions}
\item Masse molaire : $M=12\times12{,}0+22\times1{,}0+11\times16{,}0=\qty{342,0}{g.mol^{-1}}$.
\item Quantité de matière : $n=\dfrac{m}{M}=\dfrac{10{,}0}{342{,}0}=\qty{2,92e-2}{mol}$.
\item Nombre de molécules : $N=n\times N_\mathrm{A}=\num{2.92e-2}\times\num{6.02e23}=\num{1.76e22}$.
\end{questions}
\end{savoirfaire}
\begin{exemple}
\qty{100}{mL} d'eau, de masse volumique \qty{1,00}{g.mL^{-1}}, ont une masse de \qty{100}{g}. Avec $M(\ce{H2O})=\qty{18,0}{g.mol^{-1}}$, ils contiennent
$n=\dfrac{100}{18{,}0}=\qty{5,56}{mol}$ d'eau, soit $\num{3.34e24}$ molécules.
\end{exemple}
\essSection{3}{Le volume molaire d'un gaz}
\begin{loi}[Volume molaire]
Dans les mêmes conditions de température et de pression, une mole de n'importe quel gaz occupe le même volume, appelé
\textbf{volume molaire} $V_\mathrm{m}$. À \qty{20}{\degreeCelsius} et \qty{1013}{hPa}, $V_\mathrm{m}=\qty{24,0}{L.mol^{-1}}$ ; à \qty{0}{\degreeCelsius} et \qty{1013}{hPa}, $V_\mathrm{m}=\qty{22,4}{L.mol^{-1}}$. Pour un
volume $V$ de gaz :
\[n=\frac{V}{V_\mathrm{m}}.\]
\end{loi}
\schema{01-volumemolaire}
\begin{exemple}
Un litre de dioxygène à \qty{20}{\degreeCelsius} et \qty{1013}{hPa} contient $n=\dfrac{1{,}0}{24{,}0}=\qty{4,17e-2}{mol}$ de dioxygène, soit une masse
$m=n\times M=\num{4.17e-2}\times32{,}0=\qty{1,33}{g}$.
\end{exemple}
\essSection{4}{Les espèces dissoutes}
\begin{definition}[Concentrations d'un soluté]
Pour un soluté de quantité de matière $n$ et de masse $m$, dissous dans un volume $V$ de solution :
\[c=\frac{n}{V}\quad(\unit{mol.L^{-1}})\qquad\qquad t=\frac{m}{V}\quad(\unit{g.L^{-1}})\qquad\qquad t=c\times M.\]
$c$ est la \textbf{concentration en quantité de matière} ; $t$, la \textbf{concentration en masse}, a été vue en seconde.
\end{definition}
\begin{exemple}
On dissout \qty{5,85}{g} de chlorure de sodium, $M=\qty{58,5}{g.mol^{-1}}$, dans une fiole jaugée de \qty{500,0}{mL}. La solution contient
$n=\dfrac{5{,}85}{58{,}5}=\qty{0,100}{mol}$ de soluté : $c=\dfrac{0{,}100}{0{,}5000}=\qty{0,200}{mol.L^{-1}}$ et $t=0{,}200\times58{,}5=\qty{11,7}{g.L^{-1}}$.
\end{exemple}
\schema{01-dissolution}
\begin{propriete}[Dilution]
Diluer une solution mère, c'est lui ajouter du solvant. La quantité de soluté prélevée se retrouve dans la solution
fille : $c_0\,V_0=c_1\,V_1$. Le \textbf{facteur de dilution} vaut $F=\dfrac{c_0}{c_1}=\dfrac{V_1}{V_0}$ : pour diluer dix fois une solution, on prélève
\qty{10,0}{mL} de solution mère que l'on introduit dans une fiole jaugée de \qty{100,0}{mL}.
\end{propriete}
\schema{01-dilution}
\essSection{5}{La composition d'un mélange}
\begin{propriete}[Quantité de matière de chaque constituant]
Pour un mélange solide ou liquide, la composition est souvent donnée en pourcentages massiques : on calcule la masse de
chaque constituant, puis sa quantité de matière avec sa propre masse molaire. Pour un mélange de gaz, les pourcentages
sont donnés en volume ; comme tous les gaz ont le même volume molaire, ils indiquent aussi la part de chaque gaz en
quantité de matière. L'air sec contient ainsi environ 78~\% de diazote et 21~\% de dioxygène.
\end{propriete}
\begin{codecote}[0.62]
\begin{pythoncol}
M_atome = {'H': 1.0, 'C': 12.0,
           'N': 14.0, 'O': 16.0}
formule = {'C': 8, 'H': 10, 'N': 4, 'O': 2}
M = 0
for element in formule:
    M = M + formule[element] * M_atome[element]
m = 0.080                        # g
n = m / M
print(M, f"{n:.3e}", f"{n * 6.02e23:.3e}")
\end{pythoncol}
\begin{consolecol}
194.0 4.124e-04 2.482e+20
\end{consolecol}
\cote
\fontsize{9.5}{12.5}\selectfont
\capacite{Capacité numérique}\par\smallskip
Le dictionnaire \code{formule} donne le nombre d'atomes de chaque élément dans une molécule de caféine. La boucle
additionne leurs masses molaires atomiques pour obtenir $M$. Le programme calcule ensuite la quantité de matière et le
nombre de molécules de caféine dans une tasse de café qui en contient \qty{80}{mg}.
\end{codecote}
\begin{attention}
Les relations s'appliquent avec des unités cohérentes : $m$ en g et $M$ en \unit{g.mol^{-1}} ; $V$ en L et $c$ en \unit{mol.L^{-1}}. Le volume molaire
ne concerne que les gaz. Ne pas confondre la concentration en masse $t$ d'un soluté et la masse volumique $\rho$ d'une
solution : elles ont la même unité, mais l'une concerne le soluté, l'autre toute la solution.
\end{attention}
\begin{experience}[Mesurer la quantité de gaz libérée par un comprimé effervescent]
On place un comprimé effervescent et de l'eau dans un flacon relié à une éprouvette graduée retournée sur une cuve d'eau :
le dioxyde de carbone dégagé chasse l'eau de l'éprouvette, et l'on lit son volume $V$, d'où $n=V/V_\mathrm{m}$. On peut aussi
peser le flacon ouvert avant et après la réaction : la perte de masse est celle du gaz qui s'est échappé, d'où $n=m/M$. Les
deux méthodes doivent donner des valeurs voisines.
\end{experience}
\begin{rediger}
Écrire la relation utilisée, puis l'expression littérale de la grandeur cherchée. Convertir les données dans les unités
de la relation avant de remplacer. Donner le résultat avec son unité et autant de chiffres significatifs que la donnée la
moins précise.
\end{rediger}
\exercices

\rubrique{Découverte}
\begin{blocLivre}{ex-1-1}
\begin{exercice}[1]{Calculer des masses molaires}
\begin{donnees}
Masses molaires atomiques (\unit{g.mol^{-1}}) : H 1,0 ; C 12,0 ; N 14,0 ; O 16,0 ; S 32,1.
\end{donnees}
\begin{questions}
\item Calculer la masse molaire de la caféine, $\ce{C8H10N4O2}$, et celle du glucose, $\ce{C6H12O6}$.
\item Calculer la masse molaire de l'aspirine, $\ce{C9H8O4}$. Deux espèces différentes peuvent-elles avoir la même masse
molaire ?
\item Calculer la masse molaire de l'ion sulfate, $\ce{SO4^{2-}}$. Pourquoi ne tient-on pas compte des deux électrons gagnés ?
\item Calculer la masse molaire de l'ammoniac, $\ce{NH3}$.
\end{questions}
\end{exercice}
\end{blocLivre}
\begin{blocLivre}{ex-1-2}
\begin{exercice}[1]{Un diamant et une bague}
Un diamant d'un carat a une masse de \qty{0,200}{g} ; il est constitué d'atomes de carbone. Une bague en or pur a une
masse de \qty{5,0}{g}.
\begin{donnees}
$M(\ce{C})=\qty{12,0}{g.mol^{-1}}$ ; $M(\ce{Au})=\qty{197,0}{g.mol^{-1}}$ ; $N_\mathrm{A}=\qty{6,02e23}{mol^{-1}}$.
\end{donnees}
\begin{questions}
\item Calculer la quantité de matière de carbone dans le diamant, puis le nombre d'atomes de carbone.
\item Mêmes questions pour les atomes d'or de la bague.
\item Lequel des deux objets contient le plus d'atomes ? Pourquoi la comparaison des masses ne suffit-elle pas ?
\item Calculer la masse d'un atome de carbone.
\end{questions}
\end{exercice}
\end{blocLivre}
\begin{blocLivre}{ex-1-3}
\begin{exercice}[1]{Un liquide : l'éthanol}
On prélève \qty{25,0}{mL} d'éthanol, $\ce{C2H6O}$, liquide de masse volumique \qty{0,789}{g.mL^{-1}}.
\begin{donnees}
H 1,0 ; C 12,0 ; O 16,0 (\unit{g.mol^{-1}}) ; $N_\mathrm{A}=\qty{6,02e23}{mol^{-1}}$.
\end{donnees}
\begin{questions}
\item Calculer la masse d'éthanol prélevée.
\item Calculer la masse molaire de l'éthanol, puis la quantité de matière prélevée.
\item Combien de molécules d'éthanol l'échantillon contient-il ?
\item Quel volume d'éthanol faut-il prélever pour disposer d'une mole ?
\end{questions}
\end{exercice}
\end{blocLivre}
\begin{blocLivre}{ex-1-4}
\begin{exercice}[1]{Deux ballons de baudruche}
On gonfle deux ballons identiques, l'un avec de l'hélium, l'autre avec du dioxyde de carbone : chacun contient \qty{4,8}{L} de
gaz à \qty{20}{\degreeCelsius} et \qty{1013}{hPa}.
\begin{donnees}
$V_\mathrm{m}=\qty{24,0}{L.mol^{-1}}$ ; $M(\ce{He})=\qty{4,0}{g.mol^{-1}}$ ; $M(\ce{CO2})=\qty{44,0}{g.mol^{-1}}$ ; masse molaire moyenne de l'air : \qty{29,0}{g.mol^{-1}}.
\end{donnees}
\begin{questions}
\item Calculer la quantité de gaz contenue dans chaque ballon. Pourquoi est-elle la même ?
\item Calculer la masse de gaz contenue dans chaque ballon.
\item Calculer la masse d'air qui occuperait le même volume.
\item Expliquer pourquoi le ballon d'hélium s'envole alors que l'autre tombe.
\end{questions}
\end{exercice}
\end{blocLivre}
\begin{blocLivre}{ex-1-5}
\begin{exercice}[1]{Le sérum physiologique}
L'étiquette d'un sérum physiologique indique « chlorure de sodium 0,9~\% », c'est-à-dire \qty{0,9}{g} de chlorure de sodium
dissous pour \qty{100}{mL} de solution.
\begin{donnees}
$M(\ce{NaCl})=\qty{58,5}{g.mol^{-1}}$ ; $N_\mathrm{A}=\qty{6,02e23}{mol^{-1}}$.
\end{donnees}
\begin{questions}
\item Calculer la concentration en masse du chlorure de sodium, en \unit{g.L^{-1}}.
\item En déduire sa concentration en quantité de matière.
\item Un flacon contient \qty{500}{mL} de sérum. Calculer la masse et la quantité de chlorure de sodium dissous.
\item Le chlorure de sodium se dissout en ions $\ce{Na+}$ et $\ce{Cl-}$. Combien d'ions sodium le flacon contient-il ?
\end{questions}
\end{exercice}
\end{blocLivre}
\rubrique{Entraînement}
\begin{blocLivre}{ex-1-6}
\begin{exercice}[2]{Préparer une solution par dissolution}
On veut préparer \qty{250,0}{mL} d'une solution S$_0$ de sulfate de cuivre de concentration $c_0=\qty{0,0400}{mol.L^{-1}}$, à partir de
cristaux bleus de sulfate de cuivre pentahydraté, $\ce{CuSO4*5H2O}$.
\begin{donnees}
H 1,0 ; O 16,0 ; S 32,1 ; Cu 63,5 (\unit{g.mol^{-1}}). Une mole de cristaux contient une mole de $\ce{CuSO4}$ et cinq moles d'eau.
\end{donnees}
\begin{questions}
\item Calculer la masse molaire du sulfate de cuivre pentahydraté.
\item Calculer la quantité de matière de soluté nécessaire, puis la masse de cristaux à peser.
\item Rédiger le protocole de préparation de S$_0$ en nommant la verrerie utilisée.
\item Calculer la concentration en masse de cristaux dissous dans S$_0$.
\end{questions}
\end{exercice}
\end{blocLivre}
\begin{blocLivre}{ex-1-7}
\begin{exercice}[2]{Diluer une solution}
On dispose de la solution S$_0$ de l'exercice précédent, de concentration $c_0=\qty{0,0400}{mol.L^{-1}}$. On veut préparer \qty{100,0}{mL} d'une
solution S$_1$ de concentration $c_1=\qty{2,00e-3}{mol.L^{-1}}$.
\schema{01-ex07}
\begin{questions}
\item Calculer le facteur de dilution.
\item Calculer le volume de solution S$_0$ à prélever.
\item Choisir, en justifiant, la verrerie nécessaire parmi celle proposée.
\item La solution S$_1$ est-elle plus ou moins colorée que S$_0$ ? Calculer sa concentration en masse de cristaux dissous,
avec $M=\qty{249,6}{g.mol^{-1}}$.
\end{questions}
\end{exercice}
\end{blocLivre}
\begin{blocLivre}{ex-1-8}
\begin{exercice}[2]{Le degré alcoolique}
Le degré alcoolique d'une boisson est le volume d'éthanol, en millilitres, contenu dans \qty{100}{mL} de boisson. Un vin
titre 12,5~\% vol.
\begin{donnees}
Éthanol : $M=\qty{46,0}{g.mol^{-1}}$, masse volumique \qty{0,789}{g.mL^{-1}}. En France, un verre standard contient environ \qty{10}{g}
d'alcool pur (Santé publique France).
\end{donnees}
\begin{questions}
\item Calculer le volume, puis la masse d'éthanol contenus dans un litre de ce vin.
\item En déduire la concentration en quantité de matière d'éthanol du vin.
\item Calculer la masse d'éthanol dans un verre de \qty{10}{cL} de vin.
\item Comparer avec un demi de bière à 5~\% vol (\qty{25}{cL}) et un verre de whisky à 40~\% vol (\qty{3}{cL}). Commenter la notion de
verre standard.
\end{questions}
\end{exercice}
\end{blocLivre}
\begin{blocLivre}{ex-1-9}
\begin{exercice}[2]{Une eau minérale riche en magnésium}
L'étiquette d'une eau minérale indique : calcium \qty{549}{mg.L^{-1}} ; magnésium \qty{119}{mg.L^{-1}}. Ces éléments sont présents
sous forme d'ions $\ce{Ca^{2+}}$ et $\ce{Mg^{2+}}$.
\begin{donnees}
$M(\ce{Mg})=\qty{24,3}{g.mol^{-1}}$ ; $M(\ce{Ca})=\qty{40,1}{g.mol^{-1}}$ ; $N_\mathrm{A}=\qty{6,02e23}{mol^{-1}}$. Les apports conseillés en magnésium
sont de l'ordre de \qty{300}{mg} par jour pour une femme adulte (Anses).
\end{donnees}
\begin{questions}
\item Calculer la concentration en quantité de matière des ions magnésium, puis celle des ions calcium.
\item Combien d'ions magnésium un verre de \qty{20}{cL} contient-il ?
\item Quel volume de cette eau faudrait-il boire pour couvrir les apports conseillés en magnésium ? Commenter.
\item L'étiquette indique les concentrations en masse plutôt qu'en quantité de matière. Pourquoi, à votre avis ?
\end{questions}
\end{exercice}
\end{blocLivre}
\begin{blocLivre}{ex-1-10}
\begin{exercice}[2]{L'air d'une salle de classe}
Une salle de classe mesure \qty{8,0}{m} de long, \qty{7,0}{m} de large et \qty{3,0}{m} de haut. L'air, à \qty{20}{\degreeCelsius} et \qty{1013}{hPa}, contient
21~\% de dioxygène et 78~\% de diazote en volume.
\begin{donnees}
$V_\mathrm{m}=\qty{24,0}{L.mol^{-1}}$ ; $M(\ce{O2})=\qty{32,0}{g.mol^{-1}}$ ; $M(\ce{N2})=\qty{28,0}{g.mol^{-1}}$. Au repos, une personne rejette environ \qty{12}{L} de
dioxyde de carbone par heure.
\end{donnees}
\begin{questions}
\item Calculer le volume d'air de la salle en litres, puis la quantité de matière d'air.
\item Calculer la quantité de matière et la masse de dioxygène, puis de diazote, contenues dans la salle.
\item Trente élèves passent une heure dans la salle fermée. Calculer la quantité de dioxyde de carbone qu'ils rejettent.
\item Ce dioxyde de carbone s'ajoute à celui de l'air. Exprimer sa part dans l'air de la salle, en pourcentage puis en
parties par million (ppm). Pourquoi recommande-t-on d'aérer les salles de classe ?
\end{questions}
\end{exercice}
\end{blocLivre}
\begin{blocLivre}{ex-1-11}
\begin{exercice}[2]{Un programme pour les masses molaires}
On reprend le programme Python de l'essentiel.
\begin{questions}
\item \capacite{Capacité numérique}Que contiennent les dictionnaires \code{M\_atome} et \code{formule} ? Expliquer le rôle de la boucle \code{for}.
\item Interpréter les trois valeurs affichées.
\item Quelles lignes faut-il modifier pour traiter le glucose, $\ce{C6H12O6}$ ? Quelle masse molaire le programme affichera-t-il ?
\item Pour traiter le paracétamol, $\ce{C8H9NO2}$, un élève écrit \code{formule = \{'C': 8, 'H': 9, 'O': 2\}}. Quelle erreur a-t-il commise ?
Quelle valeur correcte faut-il obtenir ?
\end{questions}
\end{exercice}
\end{blocLivre}
\rubrique{Objectif DS}
\begin{blocLivre}{ex-1-12}
\begin{exercice}[3]{La glycémie}
\begin{docu}{1}{Un résultat d'analyse}
\begin{tabular}{@{}lll@{}}
Analyse & Résultat & Valeurs de référence \\
Glucose à jeun (glycémie) & \qty{7,8}{mmol.L^{-1}} & \qtyrange{3,9}{6,1}{mmol.L^{-1}} \\
\end{tabular}
\end{docu}
\begin{docu}{2}{Le diagnostic du diabète}
Une glycémie à jeun normale est comprise entre \qty{0,70}{g.L^{-1}} et \qty{1,10}{g.L^{-1}}. Le diabète est diagnostiqué lorsque la glycémie
à jeun est supérieure ou égale à \qty{1,26}{g.L^{-1}} lors de deux dosages successifs (Haute Autorité de santé).
\end{docu}
\begin{donnees}
Glucose : $\ce{C6H12O6}$, $M=\qty{180,0}{g.mol^{-1}}$. Volume sanguin d'un adulte : environ \qty{5,0}{L}. Un morceau de sucre a une masse
d'environ \qty{6}{g}.
\end{donnees}
\begin{questions}
\item Convertir le résultat d'analyse en \unit{g.L^{-1}}.
\item Vérifier que les valeurs de référence du laboratoire correspondent à l'intervalle donné par le document 2.
\item \initiative Ce patient est-il diabétique ? Que lui conseiller ?
\item Calculer la masse de glucose présente dans le sang de ce patient, puis l'excès par rapport à la limite haute de la
glycémie normale. Comparer à la masse d'un morceau de sucre.
\end{questions}
\end{exercice}
\end{blocLivre}
\begin{blocLivre}{ex-1-13}
\begin{exercice}[3]{Le réservoir d'une voiture à hydrogène}
\begin{docu}{1}{Une voiture à pile à combustible}
Une voiture à hydrogène de série stocke \qty{5,6}{kg} de dihydrogène, comprimé sous \qty{700}{bar} dans trois réservoirs d'un volume total
de \qty{142}{L}. Sa consommation est d'environ \qty{0,8}{kg} de dihydrogène pour \qty{100}{km}, et le constructeur annonce une autonomie
d'environ \qty{650}{km}.
\end{docu}
\begin{donnees}
$M(\ce{H2})=\qty{2,0}{g.mol^{-1}}$ ; $V_\mathrm{m}=\qty{24,0}{L.mol^{-1}}$ à \qty{20}{\degreeCelsius} et \qty{1013}{hPa} ; $N_\mathrm{A}=\qty{6,02e23}{mol^{-1}}$.
\end{donnees}
\begin{questions}
\item Calculer la quantité de matière de dihydrogène stockée, puis le nombre de molécules.
\item Quel volume ce gaz occuperait-il à \qty{20}{\degreeCelsius} et \qty{1013}{hPa} ? Combien de fois plus grand que le volume des réservoirs ?
Pourquoi faut-il comprimer le dihydrogène ?
\item Calculer le volume de dihydrogène, mesuré à \qty{20}{\degreeCelsius} et \qty{1013}{hPa}, consommé sur \qty{100}{km}.
\item \initiative Calculer l'autonomie que permettrait le réservoir plein et la comparer à celle annoncée. Proposer une
explication à l'écart.
\end{questions}
\end{exercice}
\end{blocLivre}
\begin{blocLivre}{ex-1-14}
\begin{exercice}[3]{La préparation d'une perfusion}
\begin{docu}{1}{La prescription}
Ajouter, dans une poche de \qty{500}{mL} de glucose à 5~\%, c'est-à-dire à \qty{50}{g.L^{-1}}, une ampoule de \qty{10}{mL} de chlorure de potassium à
10~\%, qui contient \qty{1,0}{g} de $\ce{KCl}$.
\end{docu}
\begin{docu}{2}{Une consigne du protocole de soins}
Dans le service, la concentration en ions potassium d'une perfusion ne doit pas dépasser \qty{40}{mmol.L^{-1}}.
\end{docu}
\begin{donnees}
$M(\ce{C6H12O6})=\qty{180,0}{g.mol^{-1}}$ ; $M(\ce{K})=\qty{39,1}{g.mol^{-1}}$ ; $M(\ce{Cl})=\qty{35,5}{g.mol^{-1}}$. Le chlorure de potassium est entièrement
dissous sous forme d'ions $\ce{K+}$ et $\ce{Cl-}$.
\end{donnees}
\begin{questions}
\item Calculer la masse et la quantité de matière de glucose contenues dans la poche.
\item Calculer la quantité de matière d'ions potassium apportée par l'ampoule, en millimoles.
\item Calculer la concentration en ions potassium dans la poche après l'ajout, puis la nouvelle concentration en masse du
glucose.
\item \initiative La prescription respecte-t-elle la consigne du service ? Combien d'ampoules pourrait-on ajouter au plus ?
\end{questions}
\end{exercice}
\end{blocLivre}
\begin{blocLivre}{ex-1-15}
\begin{exercice}[3]{De l'or 18 carats}
\begin{docu}{1}{Le titre d'un bijou}
Le titre en carats d'un alliage d'or est la masse d'or contenue dans \qty{24}{g} d'alliage : l'or pur titre 24 carats. Un bijoutier
propose une bague de \qty{4,0}{g} en or jaune 18 carats, composé en masse de 75~\% d'or, 12,5~\% d'argent et 12,5~\% de cuivre.
\end{docu}
\begin{donnees}
$M(\ce{Au})=\qty{197,0}{g.mol^{-1}}$ ; $M(\ce{Ag})=\qty{107,9}{g.mol^{-1}}$ ; $M(\ce{Cu})=\qty{63,5}{g.mol^{-1}}$.
\end{donnees}
\begin{questions}
\item Vérifier qu'un alliage contenant 75~\% d'or en masse titre bien 18 carats.
\item Calculer la masse, puis la quantité de matière de chacun des trois métaux de la bague.
\item \initiative Un client affirme que trois atomes sur quatre de sa bague sont des atomes d'or. A-t-il raison ?
\item Quel métal est le plus représenté en nombre d'atomes après l'or ? Expliquer pourquoi l'ordre n'est pas le même
qu'en masse.
\end{questions}
\end{exercice}
\end{blocLivre}
